\documentclass[12pt]{article}

% ------------------------------------------------
% PACKAGES
% ORCID NUMBER: 0000-0001-8803-3293
% ------------------------------------------------

\usepackage[utf8]{inputenc}
\usepackage[T1]{fontenc}

\usepackage{graphicx}
\usepackage{amsmath}
\usepackage{booktabs}
\usepackage{hyperref}
\usepackage{natbib}

\usepackage{geometry}
\geometry{margin=1in}

% ------------------------------------------------
% TITLE
% ------------------------------------------------

\title{Machine Learning-Based Optimization of Football Squad Composition Under Budget Constraints}

\author{Anonymous Authors}

\date{}

% ------------------------------------------------
% DOCUMENT
% ------------------------------------------------

\begin{document}

\maketitle

% ------------------------------------------------
% ABSTRACT
% ------------------------------------------------

\begin{abstract}

This study proposes a framework that combines machine learning and optimization
to support squad planning in professional football. A multilayer perceptron model
is trained to predict team performance based on individual player statistics from
the previous season. These predictions are used as the objective function in a
local search algorithm that simulates optimal player recruitment under budget
constraints. Experiments using data from major European leagues illustrate the
potential of the proposed approach to support decision-making in football squad
composition.

\end{abstract}

\textbf{Keywords:} sports analytics, football, machine learning, squad optimization, decision support

% ------------------------------------------------
\section{Introduction}
% ------------------------------------------------

The increasing availability of detailed sports data has created new opportunities
for the application of machine learning and optimization techniques in football
analytics. One of the most important strategic decisions for professional clubs
is the composition of their squads, which involves balancing player performance,
team needs, and financial constraints.

This study proposes a decision-support framework that combines machine learning
prediction with optimization techniques to simulate player recruitment under
budget limitations.

% ------------------------------------------------
\section{Related Work}
% ------------------------------------------------

Research on sports analytics has increasingly explored the use of machine learning
to predict match outcomes and team performance. Previous studies have applied
statistical models and machine learning techniques to estimate the probability
of victory, evaluate player contributions, and analyze tactical patterns.

However, fewer studies have explored the integration of predictive models with
optimization techniques to support squad planning decisions.

% ------------------------------------------------
\section{Data}
% ------------------------------------------------

This study uses player performance statistics collected from publicly available
football data sources. The dataset includes multiple seasons from major European
leagues, with player statistics linked to their respective teams and seasons.

Player statistics correspond to season \(t-1\), while the target variable is the
final number of points obtained by each team in season \(t\).

% ------------------------------------------------
\section{Performance Prediction Model}
% ------------------------------------------------

To predict team performance, a multilayer perceptron (MLP) model was developed
using player performance statistics as input variables. Prior to model training,
dimensionality reduction was performed using Principal Component Analysis (PCA)
to mitigate multicollinearity and reduce the dimensionality of the feature space.

The model was trained using data from one league and evaluated using data from
another league, avoiding information leakage between training and test samples.

\begin{figure}[h]
\centering
\includegraphics[width=0.8\linewidth]{figures/prediction_plot.png}
\caption{Relationship between predicted and actual team points.}
\end{figure}

% ------------------------------------------------
\section{Squad Optimization}
% ------------------------------------------------

The predictions generated by the machine learning model were used as an objective
function in a squad optimization framework. The objective is to maximize the
expected team performance under a predefined budget constraint.

A local search algorithm was implemented to explore the solution space through
player substitution operations, generating neighboring squad configurations.

% ------------------------------------------------
\section{Results}
% ------------------------------------------------

Experimental results demonstrate the predictive performance of the proposed model
and illustrate how the optimization procedure can generate improved squad
configurations under financial constraints.

\begin{figure}[h]
\centering
\includegraphics[width=0.8\linewidth]{figures/error_plot.png}
\caption{Prediction error as a function of actual points.}
\end{figure}

% ------------------------------------------------
\section{Discussion}
% ------------------------------------------------

The proposed framework highlights the potential of combining machine learning
and optimization techniques to support strategic decision-making in professional
football. The results suggest that predictive models can be integrated into
recruitment planning processes to explore alternative squad compositions.

% ------------------------------------------------
\section{Conclusion}
% ------------------------------------------------

This study presented a framework that integrates machine learning and local
search optimization to support squad planning decisions in football. Future
work may explore alternative predictive models and more advanced optimization
methods.

% ------------------------------------------------
\bibliographystyle{apalike}
\bibliography{references}

\end{document}